\section{Appendix - Spin Dependent Results}
\begin{figure}[htbp]
  \centering
  \includegraphics[ width=\linewidth]{figure6.pdf}
  \caption{The 90\% confidence-level upper limit (solid-black curve) using the mean of the nuclear structure functions from~\cite{PhysRevD.102.074018}, allowing for like-to-like comparison with other results. 
  The gray band indicates the uncertainty due to nuclear modeling across models from~\cite{PhysRevD.102.074018,Pirinen:2019gap,PhysRevLett.128.072502} and applies similarly to the limit from all xenon-based experiments. The dotted line shows the median $3\sigma$ observation significance from the post-fit model.
  Also shown are the WS2022 only~\cite{SR1WS}, XENONnT~\cite{XENON:2023cxc} (reinterpreted with a $-1\sigma$ power constraint), XENON1T~\cite{collaboration2019constraining}, PandaX-II~\cite{xia2019pandax}, and LUX~\cite{akerib2017limits} limits.}
  \label{fig:limitplotSDn}
\end{figure}


\begin{figure}[htbp]
  \centering
  \includegraphics[ width=\linewidth]{figure7.pdf}
  \caption{The 90\% confidence-level upper limit (solid-black curve) using the mean of the nuclear structure functions from~\cite{PhysRevD.102.074018}, allowing for like-to-like comparison with other results. 
  The gray band indicates the uncertainty due to nuclear modeling across models from~\cite{PhysRevD.102.074018,Pirinen:2019gap,PhysRevLett.128.072502} and applies similarly to the limit from all xenon-based experiments. 
  The dotted line shows the median $3\sigma$ observation significance from the post-fit model. 
  Also shown are the PICO-60~\cite{amole2019dark}, LZ WS2022 only~\cite{SR1WS}, XENONnT~\cite{XENON:2023cxc} (reinterpreted with a $-1\sigma$ power constraint), XENON1T~\cite{collaboration2019constraining}, PandaX-II~\cite{xia2019pandax}, and LUX~\cite{akerib2017limits} limits. The PICO-60 result relies on WIMP-scattering on the spin of the unpaired proton of $^{19}$F with minimal uncertainty.
  }
  \label{fig:limitplotSDp}
\end{figure}

In this Appendix we report limits on the spin-dependent WIMP-nucleon cross section using the same data selection, background modeling, and statistical inference techniques described in the main text. Results follow an identical procedure to that in Ref.~\cite{SR1WS}: signal models which describe scattering on \XeOneTwoNine\ (spin $1/2$, \SI{26.4}{\percent} natural abundance) and \XeOneThreeOne\ (spin $3/2$, \SI{21.2}{\percent} natural abundance)~\cite{XeAbundances} are constructed using the neutron- and proton-only nuclear structure functions and their uncertainties from Refs.~\cite{PhysRevD.102.074018,Pirinen:2019gap,PhysRevLett.128.072502}. The nominal limits use the mean structure functions from~\cite{PhysRevD.102.074018} and allow for a consistent comparison with previous limits from xenon-based experiments. As in Ref.~\cite{SR1WS}, an uncertainty band is calculated at each WIMP mass using the extrema of power-constrained limits obtained from the minimum and maximum nuclear structure function uncertainties provided in Refs.~\cite{PhysRevD.102.074018,Pirinen:2019gap,PhysRevLett.128.072502}.

The best-fit number of WIMP events is zero at all tested masses, from 9~GeV/$c^2$ to 100~TeV/$c^2$. The black line of Figure~\ref{fig:limitplotSDn} shows the 90\% confidence level nominal upper limit on the WIMP-neutron spin-dependent cross section as a function of mass, while the gray band shows the nuclear structure function uncertainty on the limit.  The minimum of the limit curve occurs at 46~GeV/$c^2$ at a cross section of $\sigma_{\mathrm{SD}}^{n} =$ \SI{3.7E-43}{cm\squared} (median expected limit at \SI{8.5E-43}{cm\squared} at 40~GeV/$c^2$), and a power constraint is applied at all tested masses. Figure~\ref{fig:limitplotSDp} shows the same for the WIMP-proton spin-dependent cross section. The minimum of the limit curve is at 40~GeV/$c^2$ at a cross section of $\sigma_{\mathrm{SD}}^p =$ \SI{9.8E-42}{cm\squared}, and a power constraint is applied between 9~GeV/$c^2$ and 91~GeV/$c^2$. The minimum of the median expected limit on $\sigma_{\mathrm{SD}}^p$ is \SI{2.3E-41}{cm\squared} and occurs at 40~GeV/$c^2$.
